On [DATE], the poet Saul R. Ort--who lived \& spoke upon the point of a spear; \& sought always a hillock of perspective—-passed away in a sleepy village on the south slopes of the Pyrenees. I knew him only in the final months of his life, having moved into the area, yet found him sharp-eyed and gentle-tongued, and an animating reader. 

\textit{Latino: A Pastoral} was his final work. It was still being tuned \& fiddled with to the end. It is an account of his last travels, in the Americas and in Europe, looking for a resting place.

There are some formal elements of the poem which may not be immediately obvious to the reader, and which I gathered instead from our conversations, and from notes he left behind. I would like to note some of them here, to draw the reader's attention. 

Like many poets, Ort saw his life as an oscillation between mental states which he analogized to Heaven and Hell. The plot of \textit{Latino} follows this oscillation, and the reader would do well to track to the various descents and ascents: the ``submerged'' city of Barra Funda, plunges from rocks to pools, drunken depression and herbal highs, rock-climbers racing their courses, birds swooping and perching again. Descents are also, often, voyages into some Triassic past.

To the end, Saul was in love with language, and used younger contacts to supply him with new words. Those careful networks of friendship are in evidence throughout \textit{Latino}: ``hot-girl summer,'' ``LARP,'' and ``cuffing season''; John Koenig's ``merrenness''; polyamory's ``metamour.'' Other references--Pynchon's Yoyodyne, Dick's kibble, the Beach Boys' ``feel flows'', de Chardin's ``noosphere'' —-belong squarely to the pop culture vocabulary of his own, ``silent'' generation, and the baby boomers that followed. A handful more are ancient: ``alba,'' ``ahiṃsā,'' ``scalpellino.''

What can I tell of his life? Scraps from conversation, never fact-checked. He had played brass in Chicago jazz bands, in the late fifties--including, once, with a visiting J.J. Johnson, who had ``schooled'' him ``severely.'' He spent brief years in Laguna Beach, CA and Greenwich Village, NY in the early 60s, and had watched a young Dylan--his peer--perform at a cafe once. (Cafe Frankenstein is mentioned by name in \textit{Latino}, and Dylan's vocal style compared to Neal Cassady's.)